% =====================================================================
%  Appendices — referenced from the main text. Place these AFTER
%  \appendix in your paper, or in a separate appendix file.
% =====================================================================

\appendix

% ---------------------------------------------------------------------
\section{Pseudocode: sample-once distillation}
\label{app:pseudocode}

Algorithm~\ref{alg:precompute} gives the precompute routine that
implements the sample-once-per-image release model
(\S\ref{ssec:threat}, \S\ref{ssec:impl}). It is invoked once before
the student training loop. Throughout, $\phi_T$ denotes the (frozen)
teacher encoder, $\mathrm{clip}(\cdot;\widehat{c})$ per-channel $L_2$
clipping to caps $\widehat{c}$, and $\eta\sim\mathcal{N}(0,\sigma^2)$
draws independent Gaussian noise per channel and spatial location.

\begin{algorithm}[t]
\caption{Sample-once precompute of noisy teacher bottlenecks.}
\label{alg:precompute}
\begin{algorithmic}[1]
\Require teacher $\phi_T$; training images $\{x_i\}_{i=1}^{N}$;
         per-channel caps $\widehat{c}\in\mathbb{R}^{C}$;
         per-channel noise scales $\sigma\in\mathbb{R}^{C}$
\Ensure  noisy bottleneck cache
         $\mathcal{C}:\mathrm{img\_id}\!\to\!\mathbb{R}^{C\times H\times W}$
\State $\mathcal{C} \gets \emptyset$
\For{$i = 1,\dots,N$}
  \State $z_i \gets \phi_T(x_i)$
         \Comment{forward through encoder, frozen}
  \State $z_i' \gets \mathrm{clip}(z_i;\widehat{c}) / \widehat{c}$
         \Comment{per-channel norm $\le 1$}
  \State sample $\eta_i \in \mathbb{R}^{C\times H\times W}$ with
         $\eta_i^{(c)} \!\sim\! \mathcal{N}(0,\sigma_c^{2})$
  \State $\widetilde{z}_i \gets (z_i' + \eta_i)\odot\widehat{c}$
         \Comment{denormalise}
  \State $\mathcal{C}[\mathrm{img\_id}(x_i)] \gets \widetilde{z}_i$
\EndFor
\State \Return $\mathcal{C}$
\end{algorithmic}
\end{algorithm}

\noindent During student training, the inner loop's teacher path is
replaced by a single look-up $\widetilde{z}_i \leftarrow
\mathcal{C}[\mathrm{img\_id}(x_i)]$, eliminating both per-iteration
teacher queries and per-iteration noise samples.

% ---------------------------------------------------------------------
\section{Full proof of Theorem~\ref{thm:wf}}
\label{app:proof}

We restate the optimisation problem from \eqref{eq:opt} for
convenience:
\begin{equation*}
\min_{\sigma_1,\ldots,\sigma_C>0}\; \sum_{c=1}^{C} s_c\,\sigma_c^{2}
\quad\text{s.t.}\quad
\sum_{c=1}^{C} \frac{\Delta_c^{2}}{2\,\sigma_c^{2}} \le \rho_{\mathrm{rel}}.
\end{equation*}

\paragraph{Change of variables.}
Let $u_c = 1/\sigma_c^{2}$, so $u_c > 0$ and $\sigma_c^{2} = 1/u_c$.
The problem becomes
\begin{equation}
\min_{u_1,\ldots,u_C>0}\; \sum_{c=1}^{C} \frac{s_c}{u_c}
\quad\text{s.t.}\quad
\sum_{c=1}^{C} \frac{\Delta_c^{2}}{2}\,u_c \le \rho_{\mathrm{rel}}.
\label{eq:dual}
\end{equation}
Both the objective $\sum_c s_c/u_c$ (convex in $u_c$ on the positive
orthant) and the linear constraint give a convex program with a
unique minimiser. The constraint is active at the optimum because the
objective is strictly decreasing in each $u_c$.

\paragraph{Lagrangian.}
Form the Lagrangian with multiplier $\lambda > 0$:
\begin{equation*}
\mathcal{L}(u,\lambda)
 \;=\; \sum_{c=1}^{C} \frac{s_c}{u_c}
       + \lambda\left(\sum_{c=1}^{C}\frac{\Delta_c^{2}}{2}u_c - \rho_{\mathrm{rel}}\right).
\end{equation*}
Stationarity $\partial\mathcal{L}/\partial u_c = 0$ gives
\begin{equation*}
-\frac{s_c}{u_c^{2}} + \frac{\lambda\,\Delta_c^{2}}{2} = 0
\;\;\Longrightarrow\;\;
u_c^{\star} \;=\; \sqrt{\frac{2\,s_c}{\lambda\,\Delta_c^{2}}}
\;=\; \sqrt{\frac{2}{\lambda}}\;\frac{\sqrt{s_c}}{\Delta_c}.
\end{equation*}

\paragraph{Solving for $\lambda$ via the active constraint.}
Substitute back into $\sum_c \tfrac{\Delta_c^{2}}{2}u_c^{\star} =
\rho_{\mathrm{rel}}$:
\begin{equation*}
\sum_{c=1}^{C}\frac{\Delta_c^{2}}{2}\sqrt{\frac{2}{\lambda}}\,\frac{\sqrt{s_c}}{\Delta_c}
 \;=\; \rho_{\mathrm{rel}}
\;\;\Longrightarrow\;\;
\sqrt{\frac{1}{2\lambda}}\,\sum_{c=1}^{C}\Delta_c\sqrt{s_c}
 \;=\; \rho_{\mathrm{rel}}.
\end{equation*}
Solving for $\lambda$:
\begin{equation*}
\lambda \;=\; \frac{1}{2\,\rho_{\mathrm{rel}}^{2}}
              \Bigl(\sum_{c=1}^{C}\Delta_c\sqrt{s_c}\Bigr)^{2}.
\end{equation*}

\paragraph{Optimal $\sigma_c^{\star}$.}
Substituting $\lambda$ back into $u_c^{\star}$ and using
$\sigma_c^{\star} = 1/\sqrt{u_c^{\star}}$:
\begin{equation*}
\sigma_c^{\star}
 \;=\;
 \frac{1}{(2/\lambda)^{1/4}}\,
 \frac{\Delta_c^{1/2}}{s_c^{1/4}}
 \;=\;
 \underbrace{\sqrt{\frac{1}{2\,\rho_{\mathrm{rel}}}\sum_{c=1}^{C}\Delta_c\sqrt{s_c}}}_{\displaystyle\kappa}
 \;\frac{\sqrt{\Delta_c}}{s_c^{1/4}},
\end{equation*}
which is exactly \eqref{eq:wf}.

\paragraph{Optimal objective and strict improvement.}
The optimal value of the original objective is
\begin{equation}
J_{\mathrm{WF}}^{\star}
 \;=\; \sum_{c=1}^{C} s_c\,(\sigma_c^{\star})^{2}
 \;=\; \kappa^{2}\sum_{c=1}^{C}\Delta_c\sqrt{s_c}
 \;=\; \frac{1}{2\,\rho_{\mathrm{rel}}}
       \Bigl(\sum_{c=1}^{C}\Delta_c\sqrt{s_c}\Bigr)^{2}.
\label{eq:wf-obj}
\end{equation}
For uniform $\sigma_{\mathrm{unif}}$ saturating the same budget,
$\sigma_{\mathrm{unif}}^{2} = (\sum_c \Delta_c^{2})/(2\,\rho_{\mathrm{rel}})$,
so the uniform-allocation objective is
\begin{equation}
J_{\mathrm{unif}}
 \;=\; \sigma_{\mathrm{unif}}^{2}\sum_{c}s_c
 \;=\; \frac{1}{2\,\rho_{\mathrm{rel}}}
       \Bigl(\sum_{c}\Delta_c^{2}\Bigr)\Bigl(\sum_{c}s_c\Bigr).
\label{eq:unif-obj}
\end{equation}
Comparing \eqref{eq:wf-obj} and \eqref{eq:unif-obj} via Cauchy--Schwarz
applied to the vectors $(\Delta_c)$ and $(\Delta_c\sqrt{s_c})$, or
directly via $(\sum_c \Delta_c\sqrt{s_c})^{2} \le
(\sum_c \Delta_c^{2})(\sum_c s_c)$ with equality iff $s_c$ is
constant, we obtain
\begin{equation*}
J_{\mathrm{WF}}^{\star} \;\le\; J_{\mathrm{unif}},
\end{equation*}
with strict inequality whenever the $\{s_c\}$ are not all equal.
This establishes the strict-improvement claim of Theorem~\ref{thm:wf}
and completes the proof. \qed

% ---------------------------------------------------------------------
%  BibTeX entries to add to references.bib (paste into your .bib):
%
%  @inproceedings{bun2016concentrated,
%    title={Concentrated Differential Privacy: Simplifications,
%           Extensions, and Lower Bounds},
%    author={Bun, Mark and Steinke, Thomas},
%    booktitle={Theory of Cryptography Conference (TCC)},
%    year={2016},
%  }
%
%  @inproceedings{lan2024gkd,
%    title={Gradient-Guided Knowledge Distillation for Object Detection},
%    author={Lan, Qizhen and Tian, Qing},
%    booktitle={IEEE/CVF Winter Conference on Applications of
%               Computer Vision (WACV)},
%    year={2024},
%  }
% ---------------------------------------------------------------------
